% Level 3, Week 1 -- Save Points: Git Basics (student workbook)
\chapterweek{1}{Save Points: Git Basics}{Git is a save-game system for code: every commit is a save point you can go back to.}

\section{What you will learn}
\begin{checklist}
  \item Use Terminal to move between folders.
  \item Make save points (commits) with git and give them helpful messages.
  \item Look back in time with \texttt{git log}, spot changes with \texttt{git diff}, and undo mistakes with \texttt{git restore}.
  \item Warm up your Level 2 skills: functions, lists and loops.
\end{checklist}

\section{The big idea}
In a video game, a save point lets you try something risky. If it goes wrong, you load the save and try again. \textbf{Git} gives your code save points. Each save point is called a \textbf{commit}. Every commit has a short id (like \texttt{a1b2c3d}) and a \textbf{message} that says what changed. Together the commits make a timeline: the story of your project.

\begin{center}
%% alt: A timeline of four save points with time running to the right: a1b "Start game", c3d "Add hero", e5f "Add jump" and g7h "Fix walls". A curved arrow says: broken? jump back to any save point.
\begin{tikzpicture}[
  commit/.style={circle,draw=black,line width=0.8pt,minimum size=9mm,font=\small\ttfamily},
  label/.style={font=\small,align=center,text width=2.7cm}]
  \draw[-{Stealth[length=3mm]},line width=1pt,kcgray] (-0.6,0) -- (13.4,0) node[right,font=\small,text=black]{time};
  \foreach \x/\id/\msg/\col in {0/a1b/Start game/kcsky, 3.3/c3d/Add hero/kcorange, 6.6/e5f/Add jump/kcgreen, 9.9/g7h/Fix walls/kcpurple}{
    \node[commit,fill=\col!40] at (\x,0) {\id};
    \node[label,below] at (\x,-0.6) {\enquote{\msg}};
  }
  \node[font=\small,align=center,anchor=west] at (10.5,1.3) {broken? jump back\\to any save point};
  \draw[-{Stealth[length=2.5mm]},kcvermillion,line width=1pt,dashed] (10.4,1.3) to[bend right=20] (3.5,0.55);
\end{tikzpicture}
\end{center}

Saving takes \textbf{two moves}, like taking a class photo:

\begin{center}
%% alt: How git moves your work. Your folder, the files you edit in Thonny, goes through git add to Ready to save, who is in the photo, then through git commit to History, the photo album of commits. git restore brings an old version back to your folder.
\begin{tikzpicture}[box/.style={draw=black,rounded corners=3pt,line width=0.8pt,minimum width=3.2cm,minimum height=1.5cm,align=center,font=\small}]
  \node[box,fill=kcyellow!35] (work) at (0,0) {\textbf{Your folder}\\files you edit\\in Thonny};
  \node[box,fill=kcsky!35] (stage) at (5.6,0) {\textbf{Ready to save}\\who is in\\the photo};
  \node[box,fill=kcgreen!35] (hist) at (11.2,0) {\textbf{History}\\the photo album\\of commits};
  \draw[-{Stealth[length=3mm]},line width=1pt] (work) -- node[above,font=\small\ttfamily]{git add} (stage);
  \draw[-{Stealth[length=3mm]},line width=1pt] (stage) -- node[above,font=\small\ttfamily]{git commit} (hist);
  \draw[-{Stealth[length=3mm]},line width=1pt,dashed] (hist.south) to[bend left=18] node[below,font=\small\ttfamily]{git restore} (work.south);
\end{tikzpicture}
\end{center}

\section{Terminal survival kit}
Terminal is a chat window with the computer: type a command, press \keys{Return}, read the answer.
\begin{center}
\renewcommand{\arraystretch}{1.25}
\begin{tabular}{@{}>{\ttfamily}l p{9.5cm}@{}}
\toprule
\normalfont\textbf{Type this} & \textbf{It means} \\ \midrule
pwd & Where am I? (print working directory) \\
ls & List the files in this folder. \\
cd git-quest & Go into the folder \texttt{git-quest}. \\
cd .. & Go up one folder. \\
\normalfont\keys{Tab} & Finish a long name for you. Type \texttt{cd git-q} and press \keys{Tab}. \\
\normalfont\keys{$\uparrow$} & Bring back the last command, so you do not retype it. \\
\bottomrule
\end{tabular}
\end{center}

\section{The six git commands}
\begin{center}
\renewcommand{\arraystretch}{1.25}
\begin{tabular}{@{}>{\ttfamily}l p{5.6cm} p{3.6cm}@{}}
\toprule
\normalfont\textbf{Command} & \textbf{What it does} & \textbf{Game words} \\ \midrule
git init & start a save system in this folder & new save file \\
git status & what changed since the last save? & pause menu \\
git add hero\_badge.py & choose what goes in the next save & pack your backpack \\
git commit -m "Add star" & make a save point with a message & save the game \\
git log --oneline & list the save points, newest first & load-game menu \\
git diff & show exactly what changed & spot the difference \\
git restore hero\_badge.py & throw away unsaved changes & reload last save \\
\bottomrule
\end{tabular}
\end{center}

\begin{tip}[Stuck in a strange screen?]
If you typed \texttt{git commit} without \texttt{-m}, git opens an editor called Vim. Press \keys{Esc}, type \texttt{:q!} and press \keys{Return}. Then try again with \texttt{-m "your message"}.
\end{tip}

\section{Good commit messages}
A good message finishes the sentence \enquote{This commit will \ldots}. Start with a verb, keep it short (50 letters or fewer) and say \emph{what} changed.
\begin{center}
\begin{tabular}{@{}p{5.5cm}p{5.5cm}@{}}
\toprule
\textbf{\ding{55}\ Not helpful} & \textbf{\checkmark\ Helpful} \\ \midrule
\texttt{stuff} & \texttt{Add jump to hero} \\
\texttt{update} & \texttt{Make star twice as big} \\
\texttt{fixed it} & \texttt{Fix score going below zero} \\
\bottomrule
\end{tabular}
\end{center}

\begin{exercise}[The Git Quest]
Folder: \texttt{week01/ex-git-quest}. Follow the 7 levels in \texttt{README.md}: make a folder called \texttt{git-quest} in Documents, set your nickname, \texttt{git init}, make your first commit, change the badge color and look at \texttt{git diff}, break the file on purpose and rescue it with \texttt{git restore}, then make a commit of your own.\\
\textbf{Done looks like:} \texttt{git log --oneline} shows 3 or more commits with helpful messages.\\
Write your commit ids and messages here (newest first):
\answerlines{3}
\end{exercise}

\begin{exercise}[Commit Coach]
Folder: \texttt{week01/ex-commit-coach}. Finish the functions in \texttt{starter.py}:
\begin{itemize}[itemsep=1pt]
  \item \texttt{is\_good\_message(message)} returns \texttt{False} for empty, too long, boring or one-word messages.
  \item \texttt{make\_message("add", "star shape")} returns \texttt{"Add star shape"}.
  \item \texttt{find\_secrets(lines)} returns the lines that contain words like \texttt{password} or \texttt{api\_key}.
\end{itemize}
Check with \texttt{check\_my\_work.py}. Then Run it and test 3 messages of your own.
\end{exercise}

\begin{exercise}[Log Reader (if you finish early)]
Folder: \texttt{week01/ex-log-reader}. Teach Python to read the text that \texttt{git log --oneline} prints. Each line is a commit id, a space, then the message:
\begin{pycode}
line = "5b7a913 Fix hero walking through walls"
parts = line.split(" ", 1)   # split only at the FIRST space
print(parts)   # ['5b7a913', 'Fix hero walking through walls']
\end{pycode}
Finish \texttt{parse\_log}, \texttt{count\_commits}, \texttt{find\_commits} and \texttt{oldest\_first}. Extra: paste your own \texttt{git log --oneline} into \texttt{SAMPLE\_LOG}.
\end{exercise}

\begin{tip}[Safety: history remembers everything]
Git writes your \texttt{user.name} and \texttt{user.email} into every commit, forever. Use a nickname and an \texttt{@example.com} email. Never commit passwords, real names, addresses or photos of people. Deleting a file later does not delete it from the history.
\end{tip}

\begin{vocab}
  \vocabword{git}{A program that keeps save points of a folder, so you can look back and undo.}
  \vocabword{repository (repo)}{A folder that git is tracking. Git keeps its notes in a hidden \texttt{.git} folder inside it.}
  \vocabword{commit}{A save point: a snapshot of your files plus a message and an id.}
  \vocabword{commit message}{A short label that says what changed, like \enquote{Add jump to hero}.}
  \vocabword{commit id}{The short code name of a commit, like \texttt{a1b2c3d}.}
  \vocabword{diff}{The differences between two versions: \texttt{-} lines were removed, \texttt{+} lines were added.}
  \vocabword{Terminal}{An app where you type commands for the computer.}
\end{vocab}

\begin{reflection}
Why does git need two steps (\texttt{add}, then \texttt{commit}) instead of one?
\answerlines{2}
Write the best commit message you made today, and why it is helpful:
\answerlines{2}
\end{reflection}
